\subsection[COMPLEMENTS]{COMPLEMENTS\footnote{In this number, we use results from chapters in preparation in the book on Commutative Algebra. We indicate them by the symbol ``Commutative Algebra''.}}

Lemma 4. Let K be a commutative field, V a finite dimensional vector space over K, G a finite group of automorphisms of V whose order q is invertible in K, S the symmetric algebra of V, and R the subalgebra of S consisting of the elements invariant under G. A prime ideal B of height 1 of S is ramified over \(p = B \cap R\) (Commutative Algebra) if and only if there exist a non-zero element a of V and a non-zero element f of \(V^{*}\) such that B = Sa and the pseudo-reflection \(s_{a,f}\) belongs to G. The decomposition group \(\mathcal{G}^{\mathrm{Z}}(\mathrm{B})\) is then the subgroup of elements of G leaving Ka stable, and the inertia group \(\mathcal{G}^{\mathrm{T}}(\mathrm{B})\) is the cyclic subgroup \(H_{a}\) of G consisting of the pseudo-reflections of G with vector a. The residue field S(B) of S at B is separable over the residue field R(p) of R at p, and the ramification index \(e(\mathrm{B}/\mathfrak{p})\) , equal to the coefficient of B, augmented by 1, in the divisor \(\text{div}(D_{S/R})\) of the different, is equal to \(\text{Card}(H_{a})\) .

To say that B is ramified over R means that its inertia group \(\mathcal{G}^{\mathrm{T}}(\mathrm{B})\) does not reduce to the identity, in other words that there exists \(g \neq 1\) in G such that \(g(z) \equiv z \pmod{B}\) for all \(z \in S\) . Since S is a factorial ring, B is a principal ideal Sa, and a must divide all the elements \(g(z) - z\) ( \(z \in S\) ); now, for \(z \in V\) , these elements are homogeneous of degree 1 and are all non-zero (since \(g \neq 1\) ); thus, a must be homogeneous of degree 1, in other words \(a \in V\) . Then there exists a linear form f on V such that \(g = s_{a,f}\) . Conversely, if g is a pseudo-reflection \(s_{a,f}\) different from 1, then \(g(z) \equiv z \pmod{Sa}\) for all \(z \in S\) , so g belongs to the inertia group of the prime ideal B = Sa. This proves the first assertion of the lemma and the characterizations of \(\mathcal{G}^{\mathrm{Z}}(\mathrm{B})\) and \(\mathcal{G}^{\mathrm{T}}(\mathrm{B})\) . Since q is coprime to the characteristic p of K (which is also that of S(B)), the extension S(B) of R(p) is separable (Commutative Algebra, Chap. V, §2, no. 2, Cor. of Prop. 5). The equality \(e(\mathrm{B}/\mathfrak{p}) = \mathrm{Card}(\mathrm{H}_{a})\) follows from this (Commutative Algebra). Since \(e(\mathrm{B}/\mathfrak{p})\) is coprime to p, the coefficient of B in \(\operatorname{div}(\mathrm{D}_{\mathrm{S/R}})\) is \(e(\mathrm{B}/\mathfrak{p}) - 1\) (Commutative Algebra). This proves the lemma.

Lemma 5. Let K be a commutative field, S a graded polynomial algebra over K, and R a graded subalgebra of S. Then S is a free graded R-module (Algebra, Chap. II, §11, no. 2) if and only if the following two conditions are satisfied: a) R is a graded polynomial algebra over K;

\begin{enumerate}
\def\labelenumi{\alph{enumi})}
\setcounter{enumi}{1}
\tightlist
\item
  if \((\alpha_{1},\ldots,\alpha_{s})\) is a system of generators of the K-algebra R consisting of algebraically independent homogeneous elements, this system is an S-regular sequence \(^{5}\) .
\end{enumerate}

When S is an R-module of finite type, b) is a consequence of a).

See Commutative Algebra for the proof.

THEOREM 4. Let K be a commutative field, V a finite dimensional vector space over K, S the symmetric algebra of V, G a finite group of automorphisms of V, and R the subalgebra of S consisting of the elements invariant under G. Assume that q = Card G is invertible in K. The following conditions are equivalent:

\begin{enumerate}
\def\labelenumi{(\roman{enumi})}
\item
  G is generated by pseudo-reflections;
\item
  S is a free graded R-module;
\item
  \(\mathbf{R}\) is a graded polynomial algebra over \(\mathbf{K}\) .
\end{enumerate}

The equivalence of (ii) and (iii) follows from no. 2 and Lemma 5. The implication (i) \(\Longrightarrow\) (ii) follows from Th. 1.

We show that (iii) \(\Longrightarrow\) (i). Let \(G'\) be the subgroup of \(G\) generated by the pseudo-reflections belonging to \(G\) , and let \(R'\) be the subalgebra of \(S\) consisting of the elements invariant under \(G'\) . We have \(R \subset R' \subset S\) . By Lemma 4, \(\text{div}(D_{S/R}) = \text{div}(D_{S/R'})\) , so \(\text{div}(D_{R'/R}) = 0\) . Assume then that \(R\) is a graded polynomial algebra. Since this is also the case for \(R'\) (since \(G'\) is generated by pseudo-reflections), Lemma 5 shows that the \(R\) -module \(R'\) admits a homogeneous basis \((Q_1, \ldots, Q_m)\) ; let \(q_i = \deg(Q_i)\) . Put

\[
d = \det (\mathrm{Tr} _ {R ^ {\prime} / R} (Q _ {i} Q _ {j})), \text { cf.   Algebra,   Chap.   IX, } \S 2.
\]

The fact that \(\operatorname{div}(\mathrm{D}_{\mathrm{R}^{\prime}/\mathrm{R}})\) is zero shows that \(\operatorname{div}(d)=0\) (Commutative Algebra), which means that d belongs to \(K^{*}\) . On the other hand \(\operatorname{Tr}_{\mathrm{R}^{\prime}/\mathrm{R}}(\mathrm{Q}_{i}\mathrm{Q}_{j})\) is a homogeneous element of degree \(q_{i}+q_{j}\) , and d is homogeneous of degree \(2\sum_{i}q_{i}\) . Thus, \(\sum_{i}q_{i}=0\) , i.e.~\(q_{i}=0\) for all i, which means that \(R^{\prime}=R\) , and so \(G^{\prime}=G\) by Galois theory. This proves that G is generated by pseudo-reflections. Q.E.D.

Remarks. 1) Under the assumptions of Th. 4, the product of the characteristic degrees of R is q (formula (6) and Prop. 2 (iii)), so they are coprime to the characteristic of K. This was asserted in no. 3.

\begin{enumerate}
\def\labelenumi{\arabic{enumi})}
\setcounter{enumi}{1}
\tightlist
\item
  If we do not assume that \(\operatorname{Card}(G)\) is invertible in K, we still have the implications (ii) \(\Longleftrightarrow\) (iii) (cf.~Lemma 5) and (ii) \(\Longrightarrow\) (i) (cf.~Exerc. 8); but the implication (i) \(\Longrightarrow\) (ii) is no longer true (Exerc. 9).
\end{enumerate}

PROPOSITION 6. The assumptions and notation are those of Th. 4. Let H be the set of pseudo-reflections belonging to G and distinct from 1. Assume that H generates G. For all \(g \in \mathbf{G}\) , put \(g(x) = x + f_g(x)e_g\) with \(e_g \in V\) , \(f_g \in V^*\) . Put

\[
\mathrm{D} = \prod_ {g \in \mathrm{H}} e _ {g} \in \mathrm{S}.
\]

\begin{enumerate}
\def\labelenumi{(\roman{enumi})}
\item
  The different of \(S\) over \(R\) is the principal ideal SD.
\item
  Identify \(S\) with the algebra \(K[X_1, \ldots, X_l]\) by choosing a basis \((X_1, \ldots, X_l)\) of \(V\) , and let \(P_1, \ldots, P_l\) be algebraically independent homogeneous elements generating the algebra \(R\) . Then the jacobian \(J = \det\left(\frac{\partial P_i}{\partial X_j}\right)\) is of the form \(\lambda D\) where \(\lambda \in K^*\) .
\item
  \(\sum_{i=1}^{l} (\deg(P_i) - 1) = \text{Card}(H)\) .
\item
  The set of anti-invariant elements of \(S\) is RD.
\end{enumerate}

Assertion (i) follows from Lemma 4. Assertion (ii) follows from the fact that SJ is the different of S over R. (Commutative Algebra). Assertion (iii) is obtained by writing down the fact that the homogeneous polynomials D and J are of the same degree. The proof of (iv) is then the same as that given in no. 4 (proof of Prop. 5, parts b) and d)).*

